% Appendix D: sensitivity of the value of targeting and of review to the parameters the
% main sweeps hold fixed (mean capability, the capability-resource association, the
% resource tail, pool size), to review's informativeness measured by rank correlation, to
% the signal model (review contaminated by resources, multiplicative review noise, review
% redrawn every round, the resource signal), and the gap-inequality summary. Added
% 2026-09-29 (revision 1) and extended 2026-09-29 (revision 2, M = 1000 particles). Data:
% sens_m.R, sens2_m.R, signals_m.R (harmonic model, smooth allocator, T = 2, b = 1);
% tables generated by gen_sens_tables.py; fig15-gap-gini.dat for Figure 15.
\section{Sensitivity of the value of targeting and of review}\label{app:sensitivity}

The main text varies the tail of the capability distribution, the budget, and review noise, and holds fixed the mean of the capability distribution, the association between capability and resources ($\rho = 0$), the tail of the resource distribution ($\alpha_R = 2$), the size of the applicant pool ($n = 50$), the noise of the resource signal ($\tau_R = 1$), and the form of the review signal (capability plus additive Gaussian noise, observed once). This appendix varies each of these and asks two questions: whether the ordering of fields by the value of targeting and of review survives, and what single quantity, if any, organizes the results. Every run uses two rounds, the default budget $b = 1$, the myopic funders with and without review, the complete-information benchmark, and $M = 1{,}000$ particles; 50 simulated populations per setting, 100 in the noise sweeps. The value of targeting is the complete-information optimum's gain over uniform funding as a share of uniform funding's gain over no funding; the value of review is the share of the records-only shortfall that the review signal recovers, as in Figure~\ref{fig:review-value}.

\subsection*{Mean capability}
A Pareto distribution with a fixed scale has mean $\alpha/(\alpha - 1)$, so at scale 1 the heavy-tailed field ($\alpha_K = 1.3$) has mean capability 4.3, the default field ($\alpha_K = 2$) 2, and the evenly spread field ($\alpha_K = 3.5$) 1.4: a sweep of $\alpha_K$ at fixed scale varies mean capability along with inequality, and a higher mean makes resources more binding on its own. Table~\ref{tab:sens-mean} compares that sweep with one in which the scale is set to $2(\alpha_K - 1)/\alpha_K$, so that every field has mean capability 2 (the realized means in the table fall short of 2 in the heavy-tailed field because the distribution has infinite variance; Appendix~\ref{app:simulation}). Holding the mean fixed narrows the heavy-tailed field's advantage but preserves the ordering at every noise level; Figure~\ref{fig:review-value} and every number in the main text use the fixed-mean sweep.

\begin{table}[t]
\centering
\small
\caption{The value of review by field at a fixed Pareto scale and at a fixed mean capability. Share of the no-review shortfall recovered, at review noise $\tau_K = 0.3$, $1$, and $3$; 100 simulated populations per entry; $M = 1{,}000$ particles. The mean capability shown is the realized mean over the populations.}
\label{tab:sens-mean}
\begin{tabular}{@{}lcccccccc@{}}
\toprule
& \multicolumn{4}{c}{fixed scale 1} & \multicolumn{4}{c}{fixed mean 2} \\
$\alpha_K$ & mean $K$ & $\tau_K = 0.3$ & $1$ & $3$ & mean $K$ & $\tau_K = 0.3$ & $1$ & $3$ \\
\midrule
1.3 & 3.84 & 0.96 & 0.90 & 0.70 & 1.77 & 0.94 & 0.79 & 0.50 \\
2 & 1.97 & 0.86 & 0.65 & 0.30 & 1.97 & 0.86 & 0.65 & 0.30 \\
3.5 & 1.39 & 0.57 & 0.23 & 0.03 & 1.99 & 0.63 & 0.32 & 0.06 \\
\bottomrule
\end{tabular}
\end{table}


\subsection*{Review compared at equal informativeness}
The same review noise $\tau_K$ is a more informative signal where capability is more dispersed, because the noise then moves fewer researchers past one another. Table~\ref{tab:sens-matched} therefore compares fields at equal informativeness, measured by the rank correlation between the review score and capability, interpolating within the fixed-mean sweep. At a rank correlation of 0.8, 0.55, and 0.3, the heavy-tailed field recovers 0.94, 0.80, and 0.50 of the shortfall, the default field 0.85, 0.65, and 0.32, and the evenly spread field 0.66, 0.46, and 0.18. The ordering is unchanged; the gaps between fields are smaller than at equal noise.

\begin{table}[t]
\centering
\small
\caption{The value of review by field at equal informativeness. Share of the records-only shortfall recovered when the rank correlation between the review score and capability is held at the stated value; fixed mean capability; interpolated within the noise sweep of Table~\ref{tab:sens-mean}.}
\label{tab:sens-matched}
\begin{tabular}{@{}lccc@{}}
\toprule
$\alpha_K$ & rank correlation 0.8 & 0.55 & 0.3 \\
\midrule
1.3 & 0.94 & 0.80 & 0.50 \\
2 & 0.85 & 0.65 & 0.32 \\
3.5 & 0.66 & 0.46 & 0.18 \\
\bottomrule
\end{tabular}
\end{table}


\subsection*{The association between capability and resources, the resource tail, and pool size}
Table~\ref{tab:sens-rest} varies the three remaining population parameters one at a time, at the default review noise and at Pareto scale 1. The association between capability and resources is the one that matters. As the copula parameter $\rho$ rises from $-0.5$ to $0.8$, the value of targeting falls in every field (from 0.96 to 0.69 in the heavy-tailed field, from 0.25 to 0.06 in the evenly spread one), and so does the value of review; the ordering of fields is unchanged at every $\rho$. The mechanism is the gap rule itself: when resources already sit with the capable, the gaps $cK_i - R_{i0}$ are small and even, and there is little for targeting, or for the information that targeting requires, to add. The tail of the resource distribution changes the value of targeting by at most a tenth and the value of review by less. Pool size raises the value of targeting in the heavy-tailed field, because a larger pool contains more of the rare researchers of very high capability, and raises the value of review in the evenly spread field, where a larger pool gives review more to separate; neither effect changes the ordering.

\begin{table}[t]
\centering
\small
\caption{The value of targeting (upper block) and of review (lower block) by field as the copula parameter $\rho$ coupling capability and resources, the tail of the resource distribution, and the size of the applicant pool vary. Review noise $\tau_K = 1$; Pareto scale 1; 50 simulated populations per entry; $M = 1{,}000$ particles. Default values: $\rho = 0$, $\alpha_R = 2$, $n = 50$.}
\label{tab:sens-rest}
\begin{tabular}{@{}lcccccccccccc@{}}
\toprule
& \multicolumn{4}{c}{copula parameter $\rho$} & \multicolumn{3}{c}{resource tail $\alpha_R$} & \multicolumn{4}{c}{pool size $n$} \\
$\alpha_K$ & $-0.5$ & $0$ & $0.5$ & $0.8$ & $1.3$ & $2$ & $3.5$ & $25$ & $50$ & $100$ & $200$ \\
\midrule
\multicolumn{12}{@{}l}{\emph{Value of targeting (share of uniform funding's gain)}} \\
1.3 & 0.96 & 0.91 & 0.82 & 0.69 & 0.92 & 0.91 & 0.82 & 0.81 & 0.91 & 1.06 & 1.15 \\
2 & 0.56 & 0.49 & 0.38 & 0.25 & 0.43 & 0.49 & 0.46 & 0.44 & 0.49 & 0.53 & 0.57 \\
3.5 & 0.25 & 0.20 & 0.13 & 0.06 & 0.17 & 0.20 & 0.18 & 0.20 & 0.20 & 0.21 & 0.22 \\
\addlinespace
\multicolumn{12}{@{}l}{\emph{Value of review (share of the no-review shortfall recovered)}} \\
1.3 & 0.92 & 0.90 & 0.88 & 0.85 & 0.88 & 0.90 & 0.92 & 0.87 & 0.90 & 0.94 & 0.95 \\
2 & 0.67 & 0.65 & 0.62 & 0.56 & 0.65 & 0.65 & 0.65 & 0.60 & 0.65 & 0.71 & 0.74 \\
3.5 & 0.24 & 0.22 & 0.16 & 0.01 & 0.24 & 0.22 & 0.21 & 0.14 & 0.22 & 0.27 & 0.32 \\
\bottomrule
\end{tabular}
\end{table}


\subsection*{The resource signal}
The main text holds the resource signal at noise $\tau_R = 1$. Table~\ref{tab:sens-tauR} varies it from nearly exact ($\tau_R = 0.1$) to nearly uninformative ($\tau_R = 10$) and removes it altogether. Resource information matters least where review matters most. In the heavy-tailed field the records-only shortfall moves from 0.44 to 0.48 of the complete-information gain across the whole range, and the value of review from 0.79 to 0.76: the researchers who matter there are identified by capability, and resources add little to their identification. In the evenly spread field a nearly exact resource signal cuts the records-only shortfall from 0.15 to 0.09, because once resources are known, output identifies capability; there the resource signal does a part of what review does, and the value of review falls from 0.39 with an exact resource signal to 0.24 with none. The ordering of fields is unchanged at every $\tau_R$.

\begin{table}[t]
\centering
\small
\caption{The resource signal. The records-only shortfall (upper block) and the value of review (lower block) as the noise of the resource signal $\tau_R$ varies from nearly exact to nearly uninformative, and with no resource signal; review noise $\tau_K = 1$; fixed mean capability; 50 simulated populations per entry; $M = 1{,}000$ particles.}
\label{tab:sens-tauR}
\begin{tabular}{@{}lcccccc@{}}
\toprule
& \multicolumn{5}{c}{resource-signal noise $\tau_R$} & none \\
$\alpha_K$ & $0.1$ & $0.3$ & $1$ & $3$ & $10$ & \\
\midrule
\multicolumn{7}{@{}l}{\emph{Records-only shortfall (share of the complete-information gain)}} \\
1.3 & 0.45 & 0.44 & 0.45 & 0.47 & 0.48 & 0.48 \\
2 & 0.25 & 0.25 & 0.27 & 0.29 & 0.30 & 0.30 \\
3.5 & 0.09 & 0.09 & 0.11 & 0.14 & 0.15 & 0.15 \\
\addlinespace
\multicolumn{7}{@{}l}{\emph{Value of review (share of the records-only shortfall recovered)}} \\
1.3 & 0.79 & 0.79 & 0.79 & 0.77 & 0.77 & 0.76 \\
2 & 0.68 & 0.67 & 0.65 & 0.61 & 0.59 & 0.59 \\
3.5 & 0.39 & 0.37 & 0.31 & 0.26 & 0.24 & 0.24 \\
\bottomrule
\end{tabular}
\end{table}


\subsection*{Multiplicative review noise}
The main text's review noise is additive, so a given $\tau_K$ is a smaller relative error for a researcher of high capability than for one of low capability, and the question is whether that form favors the heavy-tailed field. Table~\ref{tab:sens-mult} replaces the score by $K_i\exp(e_i)$ with $e_i$ Gaussian of standard deviation $s$, a constant relative error, with the funder knowing the form. The form does matter, in the direction the question anticipates. At equal informativeness, multiplicative noise costs the heavy-tailed field most: at a rank correlation of 0.65 between score and capability, review recovers 0.66 of the shortfall under multiplicative noise against 0.87 under additive noise in the heavy-tailed field, 0.61 against 0.75 in the default field, and 0.49 against 0.55 in the evenly spread field, because a constant relative error is a large absolute error for exactly the researchers who matter most in a heavy-tailed field. The ordering of fields holds at every $s$ and at every matched informativeness, and so does the main text's statement that noisy review retains most of its value under heavy tails: at $s = 0.6$, review retains 69 percent of its maximum value in the heavy-tailed field, 47 percent in the default field, and 22 percent in the evenly spread field. The heavy-tailed field's advantage is smaller under multiplicative noise than the main text's additive figures show, and it does not disappear.

\begin{table}[t]
\centering
\small
\caption{Multiplicative review noise. The review score is $K_i\exp(e_i)$ with $e_i$ Gaussian of standard deviation $s$, and the funder knows the form; fixed mean capability; 50 simulated populations per entry; $M = 1{,}000$ particles. The upper block gives the informativeness of the score at each $s$.}
\label{tab:sens-mult}
\begin{tabular}{@{}lccccc@{}}
\toprule
& \multicolumn{5}{c}{noise $s$} \\
$\alpha_K$ & $0.1$ & $0.3$ & $0.6$ & $1$ & $1.5$ \\
\midrule
\multicolumn{6}{@{}l}{\emph{Rank correlation between score and capability}} \\
1.3 & 0.96 & 0.83 & 0.65 & 0.48 & 0.35 \\
2 & 0.93 & 0.73 & 0.51 & 0.34 & 0.24 \\
3.5 & 0.85 & 0.55 & 0.33 & 0.21 & 0.14 \\
\addlinespace
\multicolumn{6}{@{}l}{\emph{Value of review (share of the no-review shortfall recovered)}} \\
1.3 & 0.96 & 0.88 & 0.67 & 0.43 & 0.27 \\
2 & 0.89 & 0.71 & 0.42 & 0.21 & 0.09 \\
3.5 & 0.68 & 0.39 & 0.15 & 0.05 & 0.02 \\
\bottomrule
\end{tabular}
\end{table}


\subsection*{Review observed once against review redrawn every round}
In the main text each researcher's review score is observed once and reused in every round, so a reviewer's misjudgment persists (\S\ref{sec:model}). Table~\ref{tab:sens-fresh} redraws the score every round with independent noise, so that repeated review accumulates information the way the record does. Over two rounds the second review adds 0.01 to 0.04 to the share of the shortfall recovered, most where review is noisiest, and the ordering of fields is unchanged. Over twenty rounds the difference is larger: in the default field the review funder's shortfall falls from 8 percent of the complete-information gain in round one to 3 percent by round twenty when the score is observed once (Figure~\ref{fig:records-rounds}), and to 0.8 percent when it is redrawn every round; in the heavy-tailed field, to 2 percent and to 0.7 percent. The records-only funder's shortfall in the same runs falls from 46 to 19 percent (default) and from 60 to 14 percent (heavy-tailed) either way. Repeated independent review would therefore strengthen the main text's comparison, and the once-observed score is the conservative case.

\begin{table}[t]
\centering
\small
\caption{Review observed once against review redrawn every round. Share of the records-only shortfall recovered over two rounds at review noise $\tau_K = 0.3$, $1$, and $3$; fixed mean capability; 100 (once) and 50 (redrawn) simulated populations per entry; $M = 1{,}000$ particles.}
\label{tab:sens-fresh}
\begin{tabular}{@{}lcccccc@{}}
\toprule
& \multicolumn{3}{c}{observed once (main text)} & \multicolumn{3}{c}{redrawn every round} \\
$\alpha_K$ & $\tau_K = 0.3$ & $1$ & $3$ & $0.3$ & $1$ & $3$ \\
\midrule
1.3 & 0.94 & 0.79 & 0.50 & 0.95 & 0.82 & 0.52 \\
2 & 0.86 & 0.65 & 0.30 & 0.88 & 0.68 & 0.32 \\
3.5 & 0.63 & 0.32 & 0.06 & 0.66 & 0.36 & 0.07 \\
\bottomrule
\end{tabular}
\end{table}


\subsection*{Review that partly reflects resources}
The main text's review signal observes capability with noise. Real review may also reflect the resources a researcher visibly has. Table~\ref{tab:sens-contam} replaces the signal by $K_i + \beta R_{i0} + \xi_i$, with the funder knowing $\beta$ and updating accordingly. Contamination costs least where capability is heavy-tailed: at $\beta = 1$, where the score weights resources as heavily as capability, the heavy-tailed field still recovers 0.70 to 0.76 of the shortfall, against 0.79 to 0.94 for a clean signal. It costs most where capability is evenly spread, where at $\beta = 1$ review recovers 0.17 to 0.26 of the shortfall, against 0.32 to 0.63 for a clean signal. The reason is the one behind Figure~\ref{fig:review-value}: the researchers whom review must find in a heavy-tailed field stand so far apart in capability that a resource term does not hide them.

\begin{table}[t]
\centering
\small
\caption{The value of review when the review score partly reflects resources. Share of the no-review shortfall recovered when the score is $K_i + \beta R_{i0}$ plus noise, for $\beta = 0$ (the main text's signal), $0.25$, $0.5$, and $1$; fixed mean capability; 50 simulated populations per entry (100 at $\beta = 0$); $M = 1{,}000$ particles.}
\label{tab:sens-contam}
\begin{tabular}{@{}lcccccccc@{}}
\toprule
& \multicolumn{4}{c}{$\tau_K = 0.3$} & \multicolumn{4}{c}{$\tau_K = 1$} \\
$\alpha_K$ & $\beta = 0$ & $0.25$ & $0.5$ & $1$ & $\beta = 0$ & $0.25$ & $0.5$ & $1$ \\
\midrule
1.3 & 0.94 & 0.92 & 0.88 & 0.76 & 0.79 & 0.77 & 0.75 & 0.70 \\
2 & 0.86 & 0.82 & 0.75 & 0.55 & 0.65 & 0.62 & 0.59 & 0.52 \\
3.5 & 0.63 & 0.55 & 0.42 & 0.26 & 0.32 & 0.27 & 0.23 & 0.17 \\
\bottomrule
\end{tabular}
\end{table}


\subsection*{The inequality of the gaps organizes the results}
The runs above vary the population parameters in five ways. One quantity summarizes what they do to the value of targeting: how unequally the optimal grants, the gaps $\max(cK_i - R_{i0}, 0)$ at the population's budget, are distributed across the field. Across all 159 population settings in this appendix (mean capability, association, resource tail, pool size, and the two noise sweeps), the logarithm of the value of targeting is a nearly linear function of the Gini coefficient of the optimal grants, with $R^2 = 0.96$ (Figure~\ref{fig:gap-gini}); adding the ratio of mean capability to mean resources raises $R^2$ to $0.99$. Settings that differ in the capability tail, the association between capability and resources, the resource tail, and pool size, but agree on the inequality of the gaps, have the same value of targeting. The value of review depends on the same quantity together with the signal's informativeness: the log-odds of the share recovered, regressed on the Gini coefficient of the optimal grants and the log-odds of the rank correlation between score and capability, has $R^2 = 0.90$, and adding the capability tail as a separate variable raises it by less than $0.01$. Capability inequality is the main driver of gap inequality in every sweep we ran, which is why the main text can speak of it directly; the association between capability and resources is the second driver, and it works through the same channel. Two cautions apply to using this quantity. The gaps depend on the budget, through $c$, as well as on the population, so the inequality of the gaps is a property of a field and a funder together; and it is a latent quantity, computed from capabilities and resources the funder does not observe, so it organizes the model rather than serving as a statistic a funder can estimate directly. The main text's Discussion lists the observable quantities that stand in for it.

% Figure (Appendix E): the value of targeting against the Gini coefficient of the optimal
% grants, across the 159 settings of Appendix E (capability tail, capability-resource
% correlation, resource tail, pool size, mean normalization). Marker by capability tail.
% Data: fig15-gap-gini.dat (gini, targ, share, skcor, a) from sens.R.
\begin{figure}[t]
\centering
\pgfplotsset{discard if not/.style 2 args={x filter/.append code={\edef\tempa{\thisrow{#1}}\edef\tempb{#2}\ifx\tempa\tempb\else\def\pgfmathresult{inf}\fi}}}
\begin{tikzpicture}
\begin{axis}[
  width=0.78\textwidth, height=6.0cm,
  axis lines*=left,
  ymode=log,
  xmin=0.25, xmax=0.95, ymin=0.04, ymax=2,
  xtick={0.3, 0.5, 0.7, 0.9},
  ytick={0.05, 0.1, 0.2, 0.5, 1, 2},
  yticklabels={0.05, 0.1, 0.2, 0.5, 1, 2},
  xlabel={Gini coefficient of the optimal grants},
  ylabel={value of targeting (share of uniform funding's gain)},
  ylabel style={font=\small}, xlabel style={font=\small},
  tick label style={font=\small},
  axis line style={line width=0.4pt},
  tick style={line width=0.4pt, black},
  clip=false,
]
\addplot[only marks, mark=*, mark size=1.4pt, black, discard if not={a}{1.3}] table[x=gini, y=targ] {fig15-gap-gini.dat};
\addplot[only marks, mark=o, mark size=1.5pt, black, line width=0.5pt, discard if not={a}{2}] table[x=gini, y=targ] {fig15-gap-gini.dat};
\addplot[only marks, mark=triangle*, mark size=1.6pt, gray, discard if not={a}{3.5}] table[x=gini, y=targ] {fig15-gap-gini.dat};
\node[font=\small, anchor=west] at (axis cs:0.27,1.55) {$\alpha_K = 1.3$ (filled)};
\node[font=\small, anchor=west] at (axis cs:0.27,1.05) {$\alpha_K = 2$ (open)};
\node[font=\small, anchor=west] at (axis cs:0.27,0.72) {$\alpha_K = 3.5$ (triangles)};
\end{axis}
\end{tikzpicture}
\caption{The value of targeting against the inequality of the optimal grants. Each point is one of the 159 settings of Appendix~\ref{app:sensitivity}: the value of targeting (logarithmic scale) against the Gini coefficient of the optimal grants at that setting's budget. The settings vary the capability tail (marker), the association between capability and resources, the resource tail, the size of the applicant pool, and the scale of capability. The logarithm of the value of targeting is linear in the Gini coefficient with $R^2 = 0.96$.}
\label{fig:gap-gini}
\end{figure}

